% ECONOMICS_PROBLEM: {"id": "D13-434", "title": "Measuring inequality within households", "jel_code": "D13", "source_id": "OP-0458"}
% FIRST_POSTED: 2026-10-09
% ATTEMPT_STATUS: partial
% ENTRY_KIND: research_attempt
\documentclass[11pt,a4paper]{article}
\usepackage[T1]{fontenc}
\usepackage[utf8]{inputenc}
\usepackage[margin=27mm]{geometry}
\usepackage{amsmath,amssymb,amsthm,mathtools}
\usepackage[hidelinks]{hyperref}
\setlength{\parindent}{0pt}
\setlength{\parskip}{5pt}
\newtheorem{theorem}{Theorem}[section]
\newtheorem{proposition}[theorem]{Proposition}
\newtheorem{lemma}[theorem]{Lemma}
\newcommand{\E}{\mathbb E}
\newcommand{\R}{\mathbb R}
\newcommand{\Prb}{\mathbb P}
\newcommand{\ind}{\mathbf1}
\DeclareMathOperator{\Var}{Var}
\DeclareMathOperator{\Cov}{Cov}
\DeclareMathOperator{\KL}{KL}
\title{Sharp adult-poverty bounds with private sharing and public services}
\author{GPT 6 Astra Ultra}
\date{9 October 2026}
\begin{document}\maketitle
\textbf{Problem ID:} \texttt{D13-434}. \textbf{Status:} Partial. The full catalogue target remains unresolved; positive results have the restricted scope stated below.

\section{A money metric is a required primitive}
The canonical target is the fraction of two adults whose individual money metric falls below a common threshold, not household poverty and not children's poverty. Household resources alone generally do not reveal each adult's welfare. This attempt specifies a simple common reference technology, derives exact sharing-rule bounds including threshold equalities, and shows that Pareto efficiency by itself does not resolve allocation.

For each household, let $S\ge0$ be total resources available for the two adults' private numeraire consumption after known spending on joint services and children. Let $g\ge0$ be an observed jointly enjoyed service amount. Adult welfare is $u_j(c_j,g)=c_j+\gamma_jg$, where $\gamma_j\ge0$ are known service valuations. The common reference environment supplies only a private numeraire $q$ at unit price with reference utility $u_j^{\rm ref}(q)=q$. Its expenditure function is $e_j(1,u)=u$ for the nonnegative utility levels here, so
\[
 m_1=c_1+\gamma_1g,\qquad m_2=c_2+\gamma_2g.
\]
This is an explicit reference convention, not an inference from ordinal preferences. The joint service may benefit both adults without twice consuming the underlying resource; money metrics are welfare equivalents, not a resource accounting identity.

Observed assignable expenditures give lower bounds $c_1\ge x_1$, $c_2\ge x_2$. More general reliable measurements may imply a closed feasible interval $c_1\in I=[l,u]\subseteq[0,S]$, with $c_2=S-c_1$. Assume $l\le u$ and all quantities measurable. The baseline expenditure restriction gives $l=x_1,u=S-x_2$ when $x_1+x_2\le S$; an empty interval rejects the measurement and accounting restrictions.

For poverty threshold $z>0$, set
\[
 a=z-\gamma_1g,\qquad b=S+\gamma_2g-z.
\]
The number of poor adults at allocation $c=c_1$ is exactly
\[
 n(c)=\ind\{c<a\}+\ind\{c>b\}. \tag{1}
\]
Strict inequalities are essential: equality to the poverty threshold is nonpoor in the catalogue definition.

\begin{theorem}[Exact conditional bounds]
The smallest feasible poor-adult count is
\[
 n_{\min}=\begin{cases}
 0,&\max(l,a)\le\min(u,b),\\
 1,&\text{otherwise, and }u\ge a\text{ or }l\le b,\\
 2,&\text{otherwise}.
 \end{cases} \tag{2}
\]
The largest feasible count is
\[
 n_{\max}=\begin{cases}
 2,&b<a,\quad u>b,\quad l<a,\\
 1,&\text{otherwise, and }l<a\text{ or }u>b,\\
 0,&\text{otherwise}.
 \end{cases} \tag{3}
\]
Every stated extremum is attained, including singleton feasible intervals.
\end{theorem}
\begin{proof}
Neither adult is poor precisely when $a\le c\le b$. This intersects $[l,u]$ exactly under the first condition in (2). If that intersection is empty, at least one adult is poor at every allocation. At most one is poor somewhere exactly when some $c\in[l,u]$ obeys $c\ge a$ or $c\le b$, equivalent to $u\ge a$ or $l\le b$. Otherwise both are always poor, completing (2).

Both are poor precisely when $b<c<a$. This open interval intersects the closed nonempty interval $[l,u]$ exactly when $b<a,u>b,l<a$. These inequalities also correctly handle $l=u$ strictly inside $(b,a)$. If no allocation has two poor adults, at least one is poor somewhere exactly when some feasible $c<a$ or $c>b$, equivalent to $l<a$ or $u>b$. Otherwise the count is always zero. For open-interval intersections, an interior midpoint or a feasible singleton attains the count. Closed conditions are attained by endpoints, proving the final assertion.
\end{proof}

\section{Population sharpness and useful boundary checks}
Let the household population law be known, and suppose no restrictions couple latent allocations across households. Then the sharp bounds on the catalogue's adult poverty rate are
\[
 P\in\left[\E[n_{\min}/2],\ \E[n_{\max}/2]\right]. \tag{4}
\]
To prove sharpness, select a minimizing or maximizing feasible allocation household by household using the endpoint and midpoint constructions in the proof. Those selections are measurable functions of $(S,g,l,u,\gamma_1,\gamma_2)$. They preserve every observation and attain the integrated endpoints. An independent common-probability randomization between the two conditional allocation laws attains every intermediate expected rate. Point identification occurs precisely when the two endpoint counts agree almost surely; this follows because their difference is nonnegative.

Several checks guard against boundary mistakes. With $S=1$, $g=0$, $I=[0,1]$, and $z=0.6$, one has $a=0.6,b=0.4$. Both-poor allocations exist, while zero-poor allocations do not, giving household poverty fractions from $1/2$ to one. With $z=0.5$, $a=b=0.5$ and equal sharing makes neither adult poor. All unequal allocations have exactly one poor adult, so the bounds become $[0,1/2]$. Changing $<$ to $\le$ would incorrectly reverse the equal-sharing boundary conclusion.

For a public-service example set $S=0.8,g=0.2$, $\gamma_1=\gamma_2=1$, $z=0.6$. Again $a=b=0.4$. Equal private sharing gives both adults money metric $0.6$ and zero poverty, while an extreme allocation yields one poor adult. The public service changes the individual thresholds without being counted as twice the household resource expenditure.

\section{Efficiency does not determine sharing}
Return to $g=0,S=1$, identical utilities $u_j=c_j$, and no informative assignable expenditures. Every full-budget allocation is Pareto efficient: increasing either person's consumption while holding the other fixed would require more than one resource unit. Moreover, every allocation maximizes the collective objective $\mu u_1+(1-\mu)u_2$ at $\mu=1/2$, which is strictly between zero and one. At $z=0.6$, allocations $(1/2,1/2)$ and $(1,0)$ are both efficient, generate identical aggregate expenditure, and yield poverty rates one and one half respectively.

Thus an efficiency restriction alone does not eliminate the ambiguity in (4). Strict concavity plus a known welfare weight might determine a unique allocation, but the weight would then supply new information. Distribution factors can identify sharing parameters only through a specified exclusion and choice model; merely labeling a variable a bargaining factor does not establish its exclusion from preferences or resources.

Measurement can narrow the interval directly. If an independent individual consumption report implies $|\widetilde c_1-c_1|\le\epsilon$, intersect $I$ with $[\widetilde c_1-\epsilon,\widetilde c_1+\epsilon]$ and reapply (2)--(4). The interval shrinks monotonically as information is added, unless it becomes empty. This is a deterministic-error illustration; noisy reports without known error envelopes require an explicit measurement distribution, and report accuracy must not be inferred from the same sharing model being validated.

The positive result is a sharp, elementary identified set for known linear service valuations and a declared reference technology. Unknown preferences, multidimensional goods, labor and time-use utility, children, endogenous public-good quantities and population sampling remain outside it. In particular, changing $\gamma_j$ or the reference service technology changes the estimand as well as its bounds. These restrictions cannot be presented as measured national individual poverty without external validation.

\begin{thebibliography}{9}
\bibitem{family} I. Alm\aa s, O. Attanasio and P. Carneiro (2023), \emph{Household decisions and intra-household distributions}, Handbook of the Economics of the Family, volume 1, chapter 3. \url{https://www.sciencedirect.com/science/chapter/handbook/pii/S2949835X23000083}. The publisher's indexed abstract and chapter text motivate preference and bargaining measurements. The exact linear-service bounds here are a restricted derivation, not a claim about the chapter's empirical estimates.
\end{thebibliography}

\end{document}
